\section{Convolution algebra and partial resolution}\label{sec:convolution}

We assume the graph has no edge loops and the stability parameters are
either $\zeta$, $\zeta^\bullet$ as in~\subsecref{subsec:Levi} or as in
\subsecref{subsec:sheaf} (and assume the graph is of an af\/f\/ine type).

\subsection{General results}

We f\/ix $\bw$ and consider {\it unions\/} of quiver varieties
$\M_{\zeta}(\bv,\bw)$, $\M_{\zeta^\bullet}(\bv,\bw)$, $\M_0(\bv,\bw)$
over va\-rious~$\bv$'s.
%
For $\M_0$ and $\M_{\zeta}$, they were introduced in
\cite[Section~2.5]{Na-qaff}:
\begin{equation*}
   \M_{0}(\bw) \defeq \bigcup_{\bv} \M_{0}(\bv,\bw),\qquad
   \M_{\zeta}(\bw) \defeq \bigsqcup_{\bv} \M_{\zeta}(\bv,\bw),
\end{equation*}
where we take the union in $\M_0(\bv,\bw)$ with respect to the closed
immersion $\M_0(V,W)\subset\M_0(V\oplus V',\bw)$ induced by the
extension by $0$ to the component $V'$. (In \cite[Section~2.5]{Na-qaff}
this was denoted by $\M_0(\infty,\bw)$.) These are inf\/inite unions of
varieties (of various dimensions), but there is no dif\/f\/iculty as we
can work on f\/initely many $\bv$'s in any of later constructions.

For $\M_{\zeta^\bullet}$ we mimic this construction. For
$\zeta^\bullet$ in \subsecref{subsec:Levi} we consider the closed
immersion $\M_{\zeta^\bullet}(V,W)\subset\M_{\zeta^\bullet}(V\oplus
V',\bw)$ with an $I^0$-graded (instead of $I$-graded) vector space
$V'$ and take the union. For $\zeta^\bullet$ in
\subsecref{subsec:sheaf} we similarly consider
$\M_{\zeta^\bullet}(V,W)\subset\M_{\zeta^\bullet}(V\oplus V',\bw)$
where $V'$ is a direct sum of various $S_i$ and $B_c$'s in
\eqref{eq:Uhbullet}.
This union is compatible with the morphism $\pi_{0,\zeta^\bullet}$ in
either cases. We have the induced morphism
$\pi_{0,\zeta^\bullet}\colon \M_{\zeta^\bullet}(\bw)\to \M_0(\bw)$. We
also have $\pi_{\zeta^\bullet,\zeta}\colon
\M_{\zeta}(\bw)\to\M_{\zeta^\bullet}(\bw)$.

Let
\(
  Z_{\zeta^\bullet,\zeta}(\bw) \defeq
  \M_{\zeta}(\bw)\times_{\M_{\zeta^\bullet}(\bw)} \M_\zeta(\bw)
\)
and similarly for $Z_{0,\zeta}(\bw)$. These are union of various
subvarieties $Z_{\zeta^\bullet,\zeta}(\bv^1,\bv^2;\bw)$
or $Z_{0,\zeta}(\bv^1,\bv^2;\bw)$
in $\M_\zeta(\bv^1,\bw)\times\M_\zeta(\bv^2,\bw)$, where the f\/iber
product is def\/ined over a space $\M_{\zeta^\bullet}(\bv,\bw)$
or $\M_0(\bv,\bw)$ with some large $\bv$ compared with $\bv^1$, $\bv^2$.
%
By \propref{prop:semismall} $Z_{0,\zeta}(\bv^1,\bv^2;\bw)$,
$Z_{\zeta^\bullet,\zeta}(\bv^1,\bv^2;\bw)$ are at most
half dimensional in $\M_\zeta(\bv^1,\bw)\times\M_\zeta(\bv^2,\bw)$.
%
Moreover $Z_{\zeta^\bullet,\zeta}(\bv^1,\bv^2;\bw)$ is a subvariety of
$Z_{0,\zeta}(\bv^1,\bv^2;\bw)$.
%
The latter space was introduced in \cite{Na-alg} as an analog of the
Steinberg variety appearing in a geometric construction of the Weyl group.

We consider the top degree Borel--Moore homology groups
$H_\topdeg(Z_{\zeta^\bullet,\zeta}(\bw))$ and
$H_\topdeg(Z_{0,\zeta}(\bw))$. More precisely the former is the subspace
\[
\prod_{\bv^1,\bv^2}' H_{\dim_\C \M_\zeta(\bv^1,\bw)\times\M_\zeta(\bv^2,\bw)}
(Z_{\zeta^\bullet,\zeta}(\bv^1,\bv^2;\bw),\Q),
\]
of the direct products consisting elements $(F_{\bv,\bv'})$ such
that{\samepage
\begin{enumerate}\itemsep=0pt
\item[1)] for f\/ixed $\bv^1$, $F_{\bv^1,\bv^2} = 0$ for all but f\/initely
many choices of $\bv^2$,
\item[2)] for f\/ixed $\bv^2$, $F_{\bv^1,\bv^2} = 0$  for all but f\/initely
many choices of $\bv^1$.
\end{enumerate}
Similarly for the latter.}

Since $\M_\zeta(\bv,\bw)$ is smooth, we have an associative algebra
structure on $H_\topdeg(Z_{\zeta^\bullet,\zeta}(\bw))$ and
$H_\topdeg(Z_{0,\zeta}(\bw))$ with the unit given by the sum of
diagonals in $\M_\zeta(\bv,\bw)\times\M_\zeta(\bv,\bw)$ for va\-rious~$\bv$.
%
We have an injective algebra homomorphism
\begin{equation}\label{eq:subalg}
    H_\topdeg(Z_{\zeta^\bullet,\zeta}(\bw)) \to
    H_\topdeg(Z_{0,\zeta}(\bw)).
\end{equation}

We can analyze irreducible representations of these algebras and the
branching rules with respect to the above homomorphism by a general
theory in \cite{CG} based on \cite{BM}. We prepare several notations
and concepts, and then state the result.

For a point $x$ in a stratum $\M_0(\bw)_{(\widehat G)}$ of
$\M_0(\bw)$ let $d_x(\bv,\bw)\! =\! \dim_\C \M_\zeta(\bv,\bw)
-\linebreak[2] \dim_\C \M_0(\bw)_{(\widehat G)}$ and
$H_\topdeg(\pi_{0,\zeta}^{-1}(x))$ be the direct sum of the cohomology
groups of the f\/ibers $\pi_{0,\zeta}^{-1}(x)\cap\M_\zeta(\bv,\bw)$ of
the degree $d_x(\bv,\bw)$.  By the convolution product
$H_\topdeg(\pi_{0,\zeta}^{-1}(x))$ is a
$H_\topdeg(Z_{0,\zeta}(\bw))$-module. The cohomology vanishes in
degree above $d_x(\bv,\bw)$ and has a basis by
$d_x(\bv,\bw)$-dimensional irreducible components in degree
$d_x(\bv,\bw)$ by \propref{prop:semismall}, and does not vanish in
$d_x(\bv,\bw)$ if the graph is of f\/inite or af\/f\/ine type by
\remref{rem:relevant}.

For a simple local system $\phi$ on $\M_0(\bw)_{(\widehat G)}$ (i.e.,
an irreducible representation of the fundamental group of
$\M_0(\bw)_{(\widehat G)}$), let $IC(\M_0(\bw)_{(\widehat G)}, \phi)$
be the corresponding intersection cohomology complex.

The fundamental group of $\M_0(\bw)_{(\widehat G)}$ acts on
$H_\topdeg(\pi_{0,\zeta}^{-1}(x))$ by monodromy. It is a permutation
of the above basis elements. We have a decomposition
\(
  H_\topdeg(\pi_{0,\zeta}^{-1}(x)) = \bigoplus \phi\otimes V_{(\widehat{G}),\phi}
\)
into the direct sum of simple local system $\phi$ tensored with the
multiplicity vector space $V_{(\widehat G),\phi}$.
%
A pair $(\M_0(\bw)_{(\widehat G)},\phi)$ of a stratum
$\M_0(\bw)_{(\widehat G)}$ and a simple local system $\phi$ on it is
{\it relevant\/} for $\pi_{0,\zeta}$ if $V_{(\widehat G),\phi}\neq 0$.
%
Then we have the decomposition theorem
\begin{equation*}
  (\pi_{0,\zeta})_*(\C_{\M_\zeta(\bw)}[\dim])
  = \bigoplus IC(\M_0(\bw)_{(\widehat G)},\phi) \otimes V_{(\widehat
    G),\phi},
\end{equation*}
where $\C_{\M_\zeta(\bw)}[\dim]$ is the direct sum
$\bigoplus_\bv \C_{\M_\zeta(\bv,\bw)}[\dim \M_\zeta(\bv,\bw)]$.
%
We also know \cite[8.9.8]{CG} that
\begin{equation}\label{eq:convalg}
  H_{\topdeg}(Z_{0,\zeta}(\bw)) \cong \bigoplus \End_\C(V_{(\widehat G),\phi}).
\end{equation}
Thus $\{ V_{(\widehat G,\phi)} \mid \text{$((\M_0)_{(\widehat
    G)},\phi)$ is relevant for $\pi_{0,\zeta}$}\}$ is the set of
isomorphism classes of irreducible representations of
$H_{\topdeg}(Z_{0,\zeta}(\bw))$.

We have similar formulas for $\pi_{\zeta^\bullet,\zeta}\colon
\M_{\zeta}(\bw)\to \M_{\zeta^\bullet}(\bw)$. We denote by $(\widehat H)$
a conjugacy class of $\GL(V)$ corresponding to a stratum of
$\M_{\zeta^\bullet}(\bw)$.
%
Similar to above we can def\/ine
$H_{\topdeg}(\pi_{\zeta^\bullet,\zeta}^{-1}(y))$, which is a
$H_\topdeg(Z_{\zeta^\bullet,\zeta}(\bw))$-module for
$y\in \M_{\zeta^\bullet}(\bw)_{(\widehat H)}$.

If we decompose the direct image of
$IC(\M_{\zeta^\bullet}(\bw)_{(\widehat H)},\psi)$ as
\begin{equation}\label{eq:bullet}
   (\pi_{0,\zeta^\bullet})_*(IC(\M_{\zeta^\bullet}(\bw)_{(\widehat
     H)},\psi))
   = \bigoplus IC(\M_0(\bw)_{(\widehat G)},\phi)\otimes
     V_{(\widehat G),\phi}^{(\widehat H),\psi}
\end{equation}
with the multiplicity vector space $V_{(\widehat G),\phi}^{(\widehat
  H),\psi}$, we have the following branching rule from the double
decomposition formula (see \cite[1.11]{BM})
\begin{equation}\label{eq:branching}
   \Res V_{(\widehat G),\phi} = \bigoplus
   V_{(\widehat G),\phi}^{(\widehat H),\psi} \otimes V_{(\widehat H),\psi},
\end{equation}
where $\Res$ is the restriction functor from
$H_{\topdeg}(Z_{0,\zeta}(\bw))$-modules to
$H_{\topdeg}(Z_{\zeta^\bullet,\zeta}(\bw))$-modules via the injective
homomorphism \eqref{eq:subalg}. Since $V_{(\widehat H),\psi}$ are
pairwise non-isomorphic, $V_{(\widehat G),\phi}^{(\widehat H),\psi}$
is the multiplicity space
\(
   \Hom_{H_{\topdeg}(Z_{\zeta^\bullet,\zeta}(\bw))}
   (V_{(\widehat H),\psi}, \Res V_{(\widehat G),\phi}).
\)


Note also that $V_{(\widehat G),\phi}^{(\widehat H),\psi}$ is the
$\phi$-isotypical component of
\(
   H^{\topdeg}(\pi_{0,\zeta^\bullet}^{-1}(x),
   IC(\M_{\zeta^\bullet}(\bw)_{(\widehat H)},\psi))[d_y]
\)
the top degree cohomology of $\pi_{0,\zeta^\bullet}^{-1}(x)$ with the
coef\/f\/icient in the shifted IC sheaf, where
$d_y = \dim \M_\zeta(\bw) - \dim \M_{\zeta^\bullet}(\bw)_{(\widehat H)}$.
(More precisely we consider the degree `$\topdeg$' and $d_y$ for each
$\bv$ separately.) See \cite[1.10]{BM}.


\subsection{The restriction to a Levi factor}\label{subsec:rest_Levi}

Let us consider the case when $\zeta$, $\zeta^\bullet$ are as in
\subsecref{subsec:Levi} in this subsection.

In \cite{Na-alg} we constructed an algebra homomorphism
\begin{equation*}
   \bU(\g) \to H_{\topdeg}(Z_{0,\zeta}(\bw)).
\end{equation*}

The homomorphism was given on generators by
\begin{equation*}
   e_i \mapsto \sum_{\bv^2} [\Pa(\bv^2,\bw)],\quad
   f_i \mapsto \sum_{\bv^2} \pm [\omega\Pa(\bv^2,\bw)], \quad
   h \mapsto \sum_{\bv} \langle h, \bw - \bv\rangle [\Delta(\bv,\bw)],
\end{equation*}
where $\Pa(\bv^2,\bw)$ is the `Hecke correspondence' parametrizing
pairs $([B^1,a^1,b^1],[B^2,a^2,b^2])$ such that $(B^1,a^1,b^1)$ is a
submodule of $(B^2,a^2,b^2)$ with the quotient isomorphic to $S_i$, a
module with $\C$ on the vertex $i$, and $0$ on the other vertices,
$\omega$ is the exchange of factors of $\M_\zeta(\bv^1,\bw)\times
\M_\zeta(\bv^2,\bw)$ and $\Delta(\bv,\bw)$ is the diagonal in
$\M_\zeta(\bv,\bw)\times \M_\zeta(\bv,\bw)$. The $\pm$-sign in the
def\/inition of $f_i$ is not important in the discussion below, so its
precise def\/inition is omitted.
%
From this def\/inition it is clear that we have an algebra homomorphism
\begin{equation*}
   \bU(\g_{I^0}+\mathfrak h) \to H_{\topdeg}(Z_{\zeta^\bullet,\zeta}(\bw)),
\end{equation*}
where $\mathfrak h$ is the Cartan subalgebra of $\g$ and $\g_{I^0}$ is
the Levi factor corresponding to the subset $I^0\subset I$.

We must be careful when we apply the result in the previous
subsection, as this homomorphism is {\it not\/} an isomorphism.

\begin{Remark}\label{rem:degree}
  Later we consider the case when $\g$ is an af\/f\/ine Lie algebra. Our
  af\/f\/ine Lie algebra, as in \cite{Kac}, contains the degree operator
  $d$. It is mapped to
\(
   \sum_\bv \langle d, \bw - \bv\rangle [\Delta(\bv,\bw)]
   = - \sum_{\bv} v_0[\Delta(\bv,\bw)],
\)
where $0$ is the special vertex of the af\/f\/ine graph as in
\subsecref{subsec:ALE}.
\end{Remark}

Let us denote by $IC(\M_{\zeta^\bullet}(\bw)_{(\widehat H)})$ the IC
sheaf corresponding to the {\it trivial\/} local system.

\begin{Theorem}\label{thm:branch}
  Let $V^\g(\lambda)$ \textup(resp.\ $V^{(\g_{I^0}+\mathfrak
    h)}(\lambda)$\textup) denote the irreducible integrable highest
  weight module of $\g$ \textup(resp.\ $\g_{I^0}+\mathfrak h$\textup)
  with the highest weight $\lambda$. Then we have
  \begin{gather*}
    (\pi_{0,\zeta^\bullet})_*(IC(\Mreg_{\zeta^\bullet}(\bv^0,\bw)))
   =
    \bigoplus_{\bv'}
    \Hom_{\g_{I^0}+\mathfrak
    h}(V^{(\g_{I^0}+\mathfrak h)}(\bw \!-\! \bv^0), V^\g(\bw\!-\!\bv'))
  \!\otimes\! IC(\Mreg_0(\bv',\bw))
  \\
  \phantom{(\pi_{0,\zeta^\bullet})_*(IC(\Mreg_{\zeta^\bullet}(\bv^0,\bw)))=}{} \oplus
  \bigoplus (\text{\rm
    $IC$ sheaves associated with {\it non\/}-regular strata}).
  \end{gather*}
\end{Theorem}

\begin{proof}
Let $\M_0(\bw)_{(\widehat G)}$ be a stratum of $\M_0(\bw)$. Take a point
$x\in \M_0(\bw)_{(\widehat G)}$. The decomposition
$H_{\topdeg}(\pi_{0,\zeta}^{-1}(x)) = \bigoplus
\phi\otimes V_{(\widehat G),\phi}$ in the previous subsection is a
decomposition of a $\g$-module.
%
Suppose $\M_0(\bw)_{(\widehat G)}$ is regular, i.e., $=
\Mreg_0(\bv',\bw)$ for some $\bv'$.
%
Then it is known that $H_{\topdeg}(\pi_{0,\zeta}^{-1}(x))$ is an
irreducible integrable representation of $\g$ with the highest weight
vector $[x] \in H_{\topdeg}(\pi_{0,\zeta}^{-1}(x)\cap
\M_\zeta(\bv',\bw))$ \cite[10.2]{Na-alg}, where
$x$ can be considered as a point in $\M_\zeta(\bv',\bw)$ as
$\pi_{0,\zeta}\colon \M_\zeta(\bv',\bw)\to \M_0(\bv',\bw)$ is an
isomorphism on the preimage of $\Mreg_0(\bv',\bw)$.
%
In particular, $V_{(\widehat G),\phi} = 0$ unless $\phi$ is the
trivial local system in this case.

Similarly if a stratum $\M_{\zeta^\bullet}(\bw)_{(\widehat H)}$ is of
the form $\Mreg_{\zeta^\bullet}(\bv^0,\bw)$ for some $\bv^0$, then
$H_{\topdeg}(\pi_{\zeta^\bullet,\zeta}^{-1}(y))$ is the irreducible
integrable representation $V^{(\g_{I^0}+\mathfrak h)}(\bw-\bv^0)$ if
$y\in \Mreg_{\zeta^\bullet}(\bv^0,\bw)$.
The highest weight vector is $[y]$ as above.

Consider a stratum $\M_{\zeta^\bullet}(\bw)_{(\widehat H)}$ which
contains $IC(\Mreg_0(\bv',\bw))$ in the decomposition
\eqref{eq:bullet}. We can write it as
$\Mreg_{\zeta^\bullet}(\bv^0,\bw)\times (\M_0)_{(\widehat H')}$ as in
\eqref{eq:Levistratum}. We consider the diagram
\eqref{eq:basechange}. Since $\kappa$ is a f\/inite birational morphism,
we have
\begin{equation*}
  IC(\overline{\Mreg_{\zeta^\bullet}(\bv^0,\bw)
     \times (\M_0)_{(\widehat H')}})
   = \kappa_*\big(
     IC(\M(\bv^0,\bw))\boxtimes IC(\overline{(\M_0)_{(\widehat H')}})
   \big),
\end{equation*}
and hence
\begin{equation*}\label{eq:nonregular}
  (\pi_{0,\zeta^\bullet})_* IC(\overline{\Mreg_{\zeta^\bullet}(\bv^0,\bw)
     \times (\M_0)_{(\widehat H')}}) = \kappa'_*\big(
   (\pi_{0,\zeta^\bullet})_*IC(\M(\bv^0,\bw))\boxtimes
   IC(\overline{(\M_0)_{(\widehat H')}})\big).
\end{equation*}
Since $\kappa'$ is also a f\/inite morphism, this contains only IC
sheaves on nonregular strata unless the second factor
$(\M_0)_{(\widehat H')}$ is trivial.
%
Hence \eqref{eq:branching} now becomes
\begin{equation}\label{eq:restriction}
   \operatorname{Res} V^{\g}(\bw - \bv')
   = \bigoplus_{\bv^0} V_{\bw-\bv'}^{\bw-\bv^0}
   \otimes V^{(\g_{I^0}+\mathfrak h)}(\bw - \bv^0),
\end{equation}
where $V_{\bw-\bv'}^{\bw-\bv^0}$ is the multiplicity of
$IC(\Mreg_0(\bv',\bw))$ in the decomposition of\linebreak[2]
$(\pi_{0,\zeta^\bullet})_*\linebreak[2](IC(\Mreg_{\zeta^\bullet}(\bv^0,
\linebreak[2]\bw)))$ in
\eqref{eq:bullet}. Thus
$V_{\bw-\bv'}^{\bw-\bv^0} =
\Hom_{\g_{I^0}+\mathfrak h}
(V^{(\g_{I^0}+\mathfrak h)}(\bw - \bv^0), V^\g(\bw-\bv'))$.
%
Note here that $V^{(\g_{I^0}+\mathfrak h)}(\bw - \bv^0)$ are
non-isomorphic for dif\/ferent $\bv^0$.
\end{proof}
